% interface=nl output=pdftex

\gebruikmodule
  [chart, abr-02]

\logo [EXAMPLE] {eXaMpLe}

\gebruikXMLfilter[mml,mmp,mmc]

\input chrt-xml

%\steluitvoerin
%  [pdftex]

\stelkorpsin
  [lbr,10pt]

\mainlanguage
  [en]

\stelwitruimtein
  [middel]

\stelblankoin
  [middel]

\steltypenin
  [blanko=middel]

%\stelpapierformaatin
%  [S6,liggend]
%  [S6,liggend]

\stellayoutin
  [breedte=midden,
   hoogte=passend,
   rugwit=12pt,
   kopwit=12pt,
   hoofd=0cm,
   voet=0cm,
   onder=8pt,
   onderafstand=16pt]

\stelkleurenin
  [status=start]

\steltypenin
  [kleur=framecolor]

\steltypein
  [kleur=linecolor]

\definecolor [focuscolor]      [r=.5,g=.9,b=.5]
\definecolor [linecolor]       [r=.9,g=.5,b=.5]
\definecolor [framecolor]      [r=.5,g=.5,b=.9]
\definecolor [backgroundcolor] [s=.50]
\definecolor [textgroundcolor] [s=.85]

\definecolor [FLOWfocuscolor]      [focuscolor]
\definecolor [FLOWlinecolor]       [linecolor]
\definecolor [FLOWframecolor]      [framecolor]
\definecolor [FLOWbackgroundcolor] [backgroundcolor]

\definieerfont[TitleFont]  [RegularBold at 60pt]
\definieerfont[AuthorFont] [RegularBold at 18pt]
\definieerfont[ContentFont][RegularBold at 15pt]

\startMPrun
  logo_type := 304 ; mpgraph := 304 ; input mp-prag ;
\stopMPrun

\stelachtergrondenin
  [pagina]
  [achtergrond=kleur,
   achtergrondkleur=backgroundcolor,
   achtergrondoffset=4pt]

\stelachtergrondenin
  [tekst][tekst]
  [achtergrond={kleur,vooruit},
   achtergrondkleur=textgroundcolor]

\stelachtergrondenin
  [status=herhaal]

\definieeroverlay[vooruit][\overlaybutton{vooruit}]

\stelinteractiein
  [kleur=backgroundcolor,
   contrastkleur=backgroundcolor,
   menu=aan,
   status=start]

%\stelinteractieschermin
%  [optie=max]

\stelsubpaginanummerin
  [status=start,
   wijze=perparagraaf]

\stelinteractiemenuin
  [onder]
  [midden=\hskip1em,
   offset=0pt,
   letter=vet,
   status=start,
   kleur=textgroundcolor,
   contrastkleur=textgroundcolor]

\startuniqueMPgraphic{frame}
  draw unitsquare
    xscaled \overlaywidth yscaled \overlayheight
    withpen pencircle scaled 1pt
    withcolor \MPcolor{\overlaylinecolor} ;
\stopuniqueMPgraphic

\definieeroverlay[frame][\uniqueMPgraphic{frame}]

\definesymbol[vorige]   [{\strut\symbool[1]}]
\definesymbol[ergens]   [{\strut\symbool[1]}]
\definesymbol[volgende] [{\strut\symbool[1]}]

\stelinteractiebalkin
  [variant=d,
   status=start,
   symbool=ja,
   kleur=textgroundcolor,
   contrastkleur=linecolor]

\startinteractiemenu[onder]
\got [SluitDocument] exit      \\
\got [eerstepagina]  begin     \\
\got [reference]     reference \\
\hfill
\txt \interactiebalk \\
\hfill
%\got [introduction]  \externalfigure[mp-xml.1]  [hoogte=8pt] \\
%\got [general]       \externalfigure[mp-xml.100][hoogte=8pt] \\
%\got [processing]    \externalfigure[mp-xml.2]  [hoogte=8pt] \\
%\got [interfaces]    \externalfigure[mp-xml.101][hoogte=8pt] \\
%\got [workflows]     \externalfigure[mp-xml.3]  [hoogte=8pt] \\
%\got [examples]      \externalfigure[mp-xml.102][hoogte=8pt] \\
%\got [reference]     \externalfigure[mp-xml.4]  [hoogte=8pt] \\
\got [introduction]  \symbol[introduction] \\
\got [general]       \symbol[general]      \\
\got [processing]    \symbol[processing]   \\
\got [interfaces]    \symbol[interfaces]   \\
\got [workflows]     \symbol[workflows]    \\
\got [examples]      \symbol[examples]     \\
\got [reference]     \symbol[reference]    \\
\stopinteractiemenu

\definefiguresymbol [introduction] [mp-xml.1]  [hoogte=8pt] 
\definefiguresymbol [general]      [mp-xml.100][hoogte=8pt] 
\definefiguresymbol [processing]   [mp-xml.2]  [hoogte=8pt] 
\definefiguresymbol [interfaces]   [mp-xml.101][hoogte=8pt] 
\definefiguresymbol [workflows]    [mp-xml.3]  [hoogte=8pt] 
\definefiguresymbol [examples]     [mp-xml.102][hoogte=8pt] 
\definefiguresymbol [reference]    [mp-xml.4]  [hoogte=8pt] 

\def\SomeLogo#1%
  {\pagina
   \naarbox
     {\externfiguur[mp-xml.#1][bfactor=max]}%
     [eerstepagina]
   \par\nointerlineskip}

\def\MainChart%
  {%\pagina
   \startregelcorrectie
   \bgroup
   \stelinteractiein[kleur=,contrastkleur=]%
   \framed
     [breedte=ruim,
      kader=overlay,
      achtergrond=frame,
      kaderkleur=linecolor]
     {\FLOWchart[convert-en]}%
   \egroup
   \stopregelcorrectie}
%   \par\nointerlineskip}

\def\SomeChart#1#2%
  {\pagina
   \bgroup
   \stelinteractiein[kleur=,contrastkleur=]%
   \setupFLOWcharts[focus=#2]%
   \naarbox
     {\framed
        [breedte=ruim,
         kader=overlay,
         achtergrond=frame,
         kaderkleur=linecolor]
        {\paginareferentie[#2]\FLOWchart[#1]}}%
     [workflows]
   \egroup
   \blanko[2*groot]
   }% \par\nointerlineskip}

\setupFLOWcharts
  [dy=.5\bodyfontsize,
   dx=2\bodyfontsize]

\def\KopCommando#1%
  {\naarbox{#1}[eerstepagina]}

\stelkopin
  [paragraaf,subparagraaf]
  [variant=paragraaf,
   nummer=nee,
   tekstcommando=\KopCommando,
   uitlijnen=links,
   pagina=,
   voor=\removelastskip,
   na=\vskip24pt,
   kleur=backgroundcolor,
   letter=\bfc]

\stelkopin
  [subparagraaf]
  [variant=paragraaf,
   nummer=nee,
   tekstcommando=\KopCommando,
   uitlijnen=links,
   voor={\blanko[2*groot]},
   na=\blanko,
   kleur=backgroundcolor,
   letter=\bfb]

\starttekst

\stelkopnummerin[paragraaf][+1]

\startstandaardopmaak

\stelinteractiein[kleur=,letter=]

\vskip24pt

\steluitlijnenin[links]

\startkleur[backgroundcolor]

\TitleFont \stelinterliniein

XML in \par \ConTeXt

\vfill

\ContentFont \stelinterliniein % better: \plaatslist[section]

\naarbox {\strut        introduction}[introduction]\par
\naarbox {\strut      general markup}[general]     \par
\naarbox {\strut    processing files}[processing]  \par
\naarbox {\strut defining interfaces}[interfaces]  \par
\naarbox {\strut     basic workflows}[workflows]   \par
\naarbox {\strut       some examples}[examples]    \par
\naarbox {\strut   command reference}[reference]   \vskip12pt

\AuthorFont \stelinterliniein

\naarbox{\strut PRAGMA ADE \blokje[breedte=1pt,hoogte=max,diepte=max] \currentdate}[laatstepagina]

\stopkleur

\vskip24pt

\stopstandaardopmaak

\SomeLogo{1} \paragraaf[introduction]{introduction}

This document is meant for those who want to use \CONTEXT\
for processing documents coded in \XML. We will introduce
some basics of \XML\ processing first. In the following
chapters we will explain the interface between \XML\ and
\CONTEXT. This document is a by||product of the internal
\PRAGMA\ project that is carried under the name \EXAMPLE,
so you may expect some examples.

Some basic knowledge of general markup as well as \CONTEXT\
does not hurt. Since everyone who read this document
probably spent some time using \CONTEXT, we hope that the
examples will ring a few bells.

We will {\em not} explain \XML, \XSLT, \TEX, \CONTEXT\ and
\PDF\ in full detail. Maybe the next chapters will inspire
the reader to download the specifications and manuals and
spend some time studying.

No, this document was not encoded in \XML, but in \TEX. The
\XML\ examples are processed as embedded \XML, so what we
keyed is what you see. We hope that you will enjoy reading
this manual.

Hans Hagen \& Ton Otten

\SomeLogo{100} \paragraaf[general]{general markup}

General markup can be defined as \citaat {coding a document
source in such a way that by reading the code, the structure
and meaning is clear}. Readability of such a source is
therefore an important prerequisite. In this chapter we will
introduce a few concepts. A more in||depth exploration of
mapping \XML\ onto \CONTEXT\ will take place in following
chapters.

\subonderwerp{coding in angle brackets}

An intentional way of coding is not new. Ever since
computers were around, users have tried to code documents
in such a way that its content could be understood by
humans and programs. As soon as the coding permits multiple
output formats, it can be classified as general markup.
General markup has nothing to do with typography, but is
focused on structure.

\starttypen
@chapter General Markup

In this chapter we will deal with ...
\stoptypen

In this example, the \type {@chapter} ensures that we can
pick up the chapter title and typeset it in a special way.
To \TEX, a more natural syntax is:

\starttypen
\chapter {General Markup}

In this chapter we will deal with ...
\stoptypen

Both examples share that they not only enable parsing the
code, but also provide a visual cue about the structure of
the document. Although some people tend to classify this
kind of coding as archaic, in principle it does not differ
too much from todays most popular coding \citeer {languages},
like \HTML\ and \XML:

\starttypen
<chapter>General Markup</chapter>

In this chapter we will deal with ...
\stoptypen

One of the purposes of coding in such a verbose way is that
we want to be able to reuse the information. Since both the
\type {@} and \type {\abc{}} way of coding is emotionally
related to applications, some people favour the \type {<xyz>}
way of dealing with source code. Although we feel that many
applications that deal with this angle bracketed code, are
just as \citeer {narrow minded} as any application, we
nevertheless follow this thread in this document.

Just adding \type {<>} to a document is not enough. A
document may look modern and up||to||date, but this kind of
neutral coding only makes sense when one is willing to look
behind the horizon of todays applications. Typesetting and
reading are moving from paper to screen and beyond, and the
more structure we add, the more prepared we are for future
reuse. But, how much structure is needed? The previous
example can be recoded as:

\starttypen
<chapter>General Markup</chapter>
<p>In this chapter we will deal with ...</p>
\stoptypen

Even these additional paragraph markers cannot be enough, so
maybe we should go further:

\starttypen
<chapter>
  <title>General Markup</title>
  <text>
    <p>In this chapter we will deal with ...</p>
  </text>
</chapter>
\stoptypen

Coding in this way, makes it easy to filter chapters, as
well as their titles, the text or even individual
paragraphs. It may be clear that in order to handle such
code properly, we also need a clear description on how to
interpret the labels. Hopefully, the labels themselves
provide some information, but it is good practice to define
them in a so called document type definition:

\starttypen
<?XML version="1.0" encoding="ISO-8859-1"?>

<!ELEMENT chapter (title,text)>
<!ELEMENT text    (p)>
<!ELEMENT title   (#PCDATA)>
<!ELEMENT p       (#PCDATA)>
\stoptypen

This means as much as: there can be \type {chapter}s, and
such a chapter can have a \type {title} and a body of \type
{text}. A text can consist of one or more paragraphs, and
since there can be many of them, we use a short element
name \type {p}.

\subonderwerp{converting \XML}

In itself, coding and defining is not enough: we want to
typeset a document, and why not use \TEX\ for doing that?
Later we will describe several ways of going from \XML\ to
for instance \PDF\ by using \TEX, but first we will
concentrate on a general transformation language, in this
case \XSLT. The next snippets define the transformation from
the previous code to \CONTEXT\ code.

\starttypen
<xsl:template match="chapter">
<xsl:apply-templates/>
</xsl:template>
\stoptypen

This is the way to say that everything between \type
{<chapter>} and \type {</chapter>} is passed along the way.

\starttypen
<xsl:template match="chapter/title">
<xsl:text>\chapter{</xsl:text>
<xsl:apply-templates/>
<xsl:text>}</xsl:text>
</xsl:template>
\stoptypen

The title is to be converted to \type {\chapter{...}}, and the
text is simply passed on.

\starttypen
<xsl:template match="chapter/text">
<xsl:apply-templates/>
</xsl:template>
\stoptypen

To make sure that \TEX\ will know where a new paragraph will
start, the \type {</p>} tags are converted into \type
{\par}'s.

\starttypen
<xsl:template match="chapter/text/p">
<xsl:apply-templates/>
<xsl:text>\par</xsl:text>
</xsl:template>
\stoptypen

A few more commands are needed to make this snippet of
\XSLT\ acceptable for programs like \type {xt}, but when
processed, we will have a proper \TEX\ file that can be
processed by \CONTEXT.

\subonderwerp{direct processing}

In many cases, the overhead involved in additional coding
is not worth the effort. Since \TEX\ and \CONTEXT\ can
manipulate text in quite advanced ways, directly coding in a
syntax more natural to \TEX, often pays off. When properly
coded, such a document can be converted to \XML\ anyway.

In \CONTEXT\ you can also mix \TEX\ and \XML\ code. You can
for instance decide to encode your formulas in \MATHML, while
the rest is just \TEX. From this you may indeed deduce that
\CONTEXT\ can handle \XML\ directly. Therefore, in many
cases, you can bypass the \XSL\ conversion stage.

\startbuffer
<math>
  <apply> <root/>
    <cn> 3 </cn>
    <apply> <times/>
      <apply> <plus/>  <ci> a </ci> <ci> b </ci> </apply>
      <apply> <minus/> <ci> c </ci> <ci> d </ci> </apply>
    </apply>
  </apply>
</math>
\stopbuffer

\typebuffer

\CONTEXT\ will convert this piece of content \MATHML\ into:

\processXMLbuffer

Although \XML\ parsing is part of the core, you explicitly
need to load the specific kind of \XML\ support you need
for your job. You define the mapping of \XML\ elements onto
\CONTEXT\ commands in your regular environments. However,
for some well standardized \DTD's \CONTEXT\ (will) provide
the corresponding \XML\ filters. Basic \MATHML\ support is
for instance part of the \CONTEXT\ distribution, although
we can imagine that you use your own implementation. The
special file prefix \type {xtag} is reserved for the
filters provided by \PRAGMA. You can load the
presentational and content \MATHML\ filters by saying:

\starttypen
\useXMLfilter[mml,mmp,mmc] % basic mathml/presentational/content
\stoptypen

\subonderwerp{getting started}

A convenient way to get familiar with \XML\ is to embed this
kind of coding into a normal \TEX\ document. In this manual,
we used buffers to accomplish this. So, given that you have
loaded the \MATHML\ filters, you can say:

\startbuffer
\startbuffer
<math>
  <apply> <eq/>
    <apply> <power/> <ci> a </ci> <ci> 2 </ci> </apply>
    <apply> <plus/>
      <apply> <power/> <ci> b </ci> <ci> 2 </ci> </apply>
      <apply> <power/> <ci> c </ci> <ci> 2 </ci> </apply>
    </apply>
  </apply>
</math>
\stopbuffer
\stopbuffer

\typebuffer \haalbuffer

You can show this buffer verbatim by saying \type {\typebuffer}
and process its content with \type {\processXMLbuffer}. As
with other buffers, you can use label them and recall them
by name.

\processXMLbuffer

Instead of buffers, you can also use:

\startbuffer
\startXMLdata
<math>
  <apply> <plus/> <ci> a </ci> <ci> b </ci> <ci> c </ci> </apply>
</math>
\stopXMLdata
\stopbuffer

\typebuffer

In print this wil become:

\haalbuffer

You can save \XML\ code in memory with:

\startbuffer
\startXMLcode[somename]
<math>
  <apply> <plus/> <ci> a </ci> <ci> b </ci> <ci> c </ci> </apply>
</math>
\stopXMLcode
\stopbuffer

\typebuffer \haalbuffer

after which you can recall it with \typ {\getXMLcode
[somename]}, for instance in the body of a macro. This
latter mechanism is meant for future communication with
other applications.

\SomeLogo{3} \paragraaf[processing]{processing files}

Apart from the previously mentioned methods of embedding
\XML\ into a document, we need a way to process a complete
\XML\ file. In doing that, we have the problem that once
\XML\ processing mode is entered, we no longer can use the
backslash to execute commands. So, we have to make sure
that we not only can enter \XML\ mode, but also quit it.

\starttypen
\defineXMLenvironment [note] {(} {)}

\enableXML just a <note>little</note> bit of XML
\stoptypen

When we process this file, \TEX\ will finally wait for you to
enter a command on the console, since it encountered the end
of a file. One solution is the following:

\starttypen
\defineXMLenvironment [note] {(} {)}

\enableXML just a <note>little</note> bit of XML

<?context-command \disableXML ?>
\stoptypen

The next solution works as well:

\starttypen
\defineXMLenvironment [note] {(} {)}

\starttext \enableXML

just a <note>little</note> bit of XML

<?context-command \stoptext ?>
\stoptypen

These methods are quite ugly, and since normally we will use
\XML\ source code that is stored in another file, in most
cases one will use one of:

\starttypen
\processXMLfile        {filename.xml}
\processXMLfilegrouped {filename.xml}
\stoptypen

This second command is probably the best one to use, as in:

\starttypen
% here go style definitions
% and XML to TEX mappings

\starttext
\processXMLfilegrouped {filename.xml}
\stoptext
\stoptypen

\SomeLogo{101} \paragraaf[interfaces]{defining interfaces}

When mapping \XML\ onto \CONTEXT\ commands, the complexity
of the definitions depends on to what extend the content is
to be manipulated. We will start with some examples of
straightforward mappings, where no tricky things are needed.

\subonderwerp{mapping elements}

Any document has a starting point and an end point.
Furthermore there will be many rather symmetrical
structuring commands.

Consider the following example:

\starttypen
<document>
  <title name="small">A very small document</title>
  <list packed="yes">
    <item>of course these commands <b>don't</b> match &context;
          commmand names</item>
    <item>and even worse, <ref name="att"> attributes will not
          be the same as in &context;</item>
  </list>
</document>
\stoptypen

For the moment we will ignore the attributes and just handle
the plain elements. The document is encapsulated in a \type
{document} element. You can use the begin- and endtag to do
initializations, but for the moment we stick to a simple
definition.

\starttypen
\defineXMLenvironment [document] \starttext    \stoptext
\stoptypen

The title element comes next. Because we need to pass the
content of this element to a context heading command, like
\type {\title}, we will define it as an argument element.

\starttypen
\defineXMLargument [title] \title
\stoptypen

This means as much as that the content is collected and
passed on to \type {\title}. As usual, \CONTEXT\ will take
care of preparing the necessary table of content entries,
if needed.

The itemized list has child elements that identify the
individual components. A suitable definition is:

\starttypen
\defineXMLenvironment [list] \startitemize \stopitemize
\defineXMLenvironment [item] \item         \par
\stoptypen

Because in \CONTEXT\ you don't have to explicitly mark the
end of an item, the last element can be defined as a
command element:

\starttypen
\defineXMLcommand [item] \item
\stoptypen

The \type {b} element can be mapped onto \CONTEXT's \type
{\bf} command. The best way to do this is the following:

\starttypen
\defineXMLgrouped [b] \bf
\stoptypen

In this case we could have said:

\starttypen
\defineXMLenvironment [b] {\bgroup\bf} {\egroup}
\stoptypen

But in other situations the previous definition is more
robust. Although a valid definition, the following one is to
be avoided, since picking up arguments takes time.

\starttypen
\defineXMLargument [b] {\groupedcommand{\bf}{}}
\stoptypen

The reference element \type {ref} is an example of an empty
element. The following invocations are equivalent:

\starttypen
<ref name="something"></ref>
<ref name="something" />
\stoptypen

For handling empty elements, we have two definition commands
available. Given that \type {\myXMLref} is defined, you can
say:

\starttypen
\defineXMLcommand  [ref] \myXMLref
\defineXMLsingular [ref] \myXMLref
\stoptypen

The command definition handles both cases, while the
\type {singular} one only handles the second (\type {../>})
alternative. This makes it possible to distinguish between
both cases, although officially there is no distinction.

Now, how should we define this \type {\myXMLref} command. It
is clear that we have to use the \type {name} attribute.
Because the definition of \type {\myXMLref} is rather
simple, we can instead directly handle the reference, as in:

\starttypen
\defineXMLsingular [ref] {\pagereference[\XMLpar{ref}{name}{}]}
\stoptypen

You can pick up a parameter with \type {\XMLpar}. Each
parameter is accessed by its element and attribute. The
third argument to \type {\XMLpar} is the default value. We
can slightly change the definition of the \type {title}
element.

\starttypen
\defineXMLargument [title] {\title[\XMLpar{title}{name}{}]}
\stoptypen

For more advanced attribute handling, there are specific
testing macros available, and to enlighten the task of using
the attributes as input for \CONTEXT\ setup commands, there
are methods to map values onto other ones.

A complication in mapping is that we have to make sure that
the values are passed in the way that macros expect them. For
instance, if we want to convert \type {packed="yes"} into
the \type {packed} directive, we can use:

\starttypen
\XMLifequalelse{itemize}{packed}{yes}{packed}{}
\stoptypen

This expands into \type {packed} or into an empty string.
Expansion can be very confusing since one never knows when
this takes place, unless one has a rather detailed
knowlegde of \CONTEXT's low level behaviour.

\starttypen
\defineXMLenvironment [itemize]
  {\startitemize[\XMLifequalelse{itemize}{packed}{yes}{packed}{}]}
  {\stopitemize}
\stoptypen

The following definition works out all right, since \type
{\startitemize} expands its argument, but in other
situations one may have to force expansion:

\starttypen
\defineXMLenvironment [itemize]
  {\expanded{\startitemize[....]}}
  {\stopitemize}
\stoptypen

This example also demonstrates that mapping \XML\ onto
\CONTEXT\ is not that hard, but sometimes needs a bit of
hard work. The good news is of course that once that job is
done, you seldom have to look into it again.

\subonderwerp{entities}

You can compare entities to \CONTEXT's named glyphs,
abbreviations, logos and alike. They are, so to say,
predefined representations of something. In the previous
example we met \type {&context;} as an example of an entity.

You can define an entity with:

\starttypen
\defineXMLentity [context] {The ConTeXt Macro Package}
\stoptypen

although, when you use \TEX, it makes more sense to say:

\starttypen
\defineXMLentity [context] \ConTeXt
\stoptypen

We will provide an entity dictionary for the more or less
standardized named glyphs, the most common logos, and
special characters. You can load these by saying:

\starttypen
\useXMLfilter[ent]
\stoptypen

As with elements, you can always overload an entity. Behind
the screens, entities become empty elements, which enables
them to move around in the \CONTEXT\ kernel and auxiliary
files without much trouble. So, don't be surprised when you
see such a disguised entity passing by on your screen.

\subonderwerp{delimited environments}

In \TEX\ there are several ways to define a command that
acts upon an argument. The most straightforward way is:

\starttypen
\def\MyCommand#1{... #1 ...}
\stoptypen

Such a command is invoked with \type {\MyCommand{Works}}
where the \type {{}} determines the begin and end of the
argument.

Sometimes, using braces is not that charming, and
delimiters are used. Many setup commands in \CONTEXT\ use
\type {[]} to delimit the parameters. One reason for this
is that it stands out in editors, especially when they
support syntax highlighting. Another reason is that
implementing optional arguments is more convenient with
non||braces.

Yet another kind of delimiting is the following:

\starttypen
\StartOfMyText
... my beautiful text ...
\StopOfMyText
\stoptypen

Such a command can be defined as:

\starttypen
\def\StartOfMyText#1\StopOfMyText{... #1 ...}
\stoptypen

Normally users of \CONTEXT\ can remain happily unaware of
how a command is implemented. There are however a few cases
where \CONTEXT\ uses delimited commands of the last
category, for instance in natural tables. Since these
macros were developed with processing \HTML\ in mind, using
them in an \XML\ environment makes sense.

\startbuffer[xml]
<table>
  <tr> <td> a </td> <td> one   </td> <td> alpha </td> </tr>
  <tr> <td> b </td> <td> two   </td> <td> beta  </td> </tr>
  <tr> <td> c </td> <td> three </td> <td> gamma </td> </tr>
</table>
\stopbuffer

\typebuffer[xml]

Here we need to map \typ {<tr> ... </tr>} onto \typ {\bTR
... \eTR} and \typ {<td> ... </td>} onto \typ {\bTD ...
\eTD}. This is handled by the pickup element
environment.

\startbuffer[def]
\defineXMLenvironment [table] \bTABLE \eTABLE
\defineXMLpickup      [tr]    \bTR    \eTR
\defineXMLpickup      [td]    \bTD    \eTD
\stopbuffer

\typebuffer[def] \haalbuffer[def]

For simple tables this solution works quite well, which is
demonstrated here.

\startregelcorrectie
\processXMLbuffer[xml]
\stopregelcorrectie

But, what happens when a user nests a table. Just take a
close look at the following table, and try to imagine what
happens when we pick up the argument of the \type {<td>}
that is followed by \type {<table>}.

\startbuffer[xml]
<table>
<tr> <td> a </td> <td> 1 </td> </tr>
<tr> <td> <table> <tr> <td> a </td> <td> 1 </td>
                       <td> b </td> <td> 2 </td>
                       <td> c </td> <td> 3 </td> </tr>
   </table> </td> <td> 2 </td> </tr>
<tr> <td> c </td> <td> 3 </td> </tr>
</table>
\stopbuffer

\typebuffer[xml]

The argument will not contain the whole table, but just
\typ {<table> <tr> <td> a} since \TEX\ will stop grabbing
the argument when it encounters the \type {</td>}. Even if
\TEX\ would have a regular expression engine like \PERL,
handling this kind of nesting would be quite complicated.

In a document coded in \TEX, putting \type {{}} around the
inner table would solve the problem, so this is what we will
do behind the screen.

\startbuffer[def]
\defineXMLenvironment [table] \bTABLE \eTABLE
\defineXMLnested      [tr]    \bTR    \eTR
\defineXMLnested      [td]    \bTD    \eTD
\stopbuffer

\typebuffer[def] \haalbuffer[def]

The nested environment analyzes its argument and keeps
picking up code until the nesting is proper. Because this
slows down processing, you should not use the nested
environment unneeded.

\startregelcorrectie
\processXMLbuffer[xml]
\stopregelcorrectie

This also works with more complicated cases, like the
following (unreadable) table:

\startbuffer[xml]
<table>
<tr> <td> a </td> <td> 1 </td> </tr>
<tr> <td> <table> <tr> <td> a </td> <td> 1 </td>
                       <td> <table>
                              <tr> <td> a </td> <td> 1 </td>
                                   <td> b </td> <td> 2 </td>
                                   <td> c </td> <td> 3 </td> </tr>
                     </table> </td> <td> 2 </td>
                       <td> c </td> <td> 3 </td> </tr>
   </table> </td> <td> 2 </td> </tr>
<tr> <td> c </td> <td> 3 </td> </tr>
</table>
\stopbuffer

\typebuffer[xml]

\startregelcorrectie
\processXMLbuffer[xml]
\stopregelcorrectie

\subonderwerp{pushing and popping}

\startbuffer
<text>

<title>Hello <tt>\World</tt>, this is &tex;!</title>

Welcome to this small XML example.

<figure label="cow">
  <caption>Some Caption</caption>
  <graphic>koe.pdf</graphic>
</figure>

<!-- Next comes some rather contextual code. The ? indicates
that we are dealing with a processing instruction. -->

<?context-command {\bf What a happy end to a small story!} ?>

</text>
\stopbuffer

\typebuffer

We will start with the easy part of the interface to this
\XML\ document.

\starttypen
\defineXMLenvironment [text]  {\starttext} {\stoptext}
\defineXMLargument    [title] {\chapter[\XMLpar{title}{label}{}]}
\defineXMLgrouped     [tt]    {\tt}
\stoptypen

In order to resolve some entities, we also say:

\starttypen
\useXMLfilter[ent]
\stoptypen

The comment is skipped, while the processing instruction is
processed right away. So, how about the graphic.

\starttypen
\defineXMLenvironment [figure]
  {\bgroup
   \defineXMLpush[caption]
   \defineXMLpush[graphic]}
  {\placefigure
     [\XMLpar{figure}{location}{here}]
     [fig:\XMLpar{figure}{label}{unknown}]
     {\XMLpop{caption}}
     {\externalfigure[\XMLpop{graphic}]}
   \egroup}
\stoptypen

Before we can process the graphic, we need to know its
caption. Therefore, we push the caption and the graphic
definition into memory and at the end pop them in the order
needed.

\subonderwerp{more on attributes}

In the previous example, the location is passed as it is.
If the \XML\ code is not yours, there is a good chance that
the values don't match those understood bij \CONTEXT. This
omission can be compensated by mapping a value, as in:

\starttypen
\mapXMLvalue{figure}{location}{thisplace}{here}
\stoptypen

When we replace the \type {\XMLpar} with \type {\XMLvar}, we
get the remapped value. This would permit calls like:

\starttypen
<figure label="cow" location="thisplace"> ... </figure>
\stoptypen

\SomeLogo{1} \paragraaf[workflows]{basic workflows}

We won't be surprised if, at this stage, you have lost track
on how to proceed in going from \XML\ to \TEX. This chapter
will describe several methods, and which way you go depends
on both your input and your wishes.

\MainChart

\SomeChart {conversion 1} {}

This manual was processed this way. The source was coded in
\TEX, but in order to demonstrate how \XML\ is processed, we
used snippets of \XML. In such cases, the necessary \XML\
support is part of the regular style file.

\starttypen
texexec --pdf yourfile
\stoptypen

\SomeChart {conversion 2} {}

The most straightforward way to process \XML\ is to let
\CONTEXT\ directly interpret the \XML\ commands. Normally you
will use a specific environment file in which the mapping as
well as the style setup is defined.

\starttypen
texexec --pdf --env=list yourfile.xml
\stoptypen

In this case, \TEXEXEC\ will make a dummy file called \type
{yourfile.tex} that loads the \XML\ file. If you simply want
to test a \MATHML\ fragment, you can say:

\starttypen
texexec --pdf --xmlfilter=mml,mmp,mmc yourfile.xml
\stoptypen

\SomeChart {conversion 3} {}

This method uses a dedicated converter that converts the
\XML\ code into a \TEX\ file where angle brackets and
ampersands are replaced by a few \TEX\ commands and where
characters that have a special meaning in \TEX\ become
entities. In most cases, the direct interpretation of \XML\
code is more suitable, but for less restrictive dialects
this is a safe way to go. Discussing this method, which is
occasionally used at \PRAGMA\ is beyond the scope of this
document.

\SomeChart {conversion 4} {}

This can be a convenient method but its usage depends on
the complexity of the mapping. Although using \XSLT\ has a
certain charm, a decent mapping file can look quite
unreadable and take some time to prepare. This method is
used when manipulations have to take place that are not
present in \CONTEXT. The resulting \TEX\ file is processed
as any normal \TEX\ file.

\SomeChart {conversion 6} {}

When you want to do a lot of manipulations, or when you want
to convert part of the \XML\ code into other code (think of
manipulating math), you can use \XSLT\ to do that before the
document source is passed on to \CONTEXT. It is to be
expected that such filters will become available.

\SomeChart {conversion 5} {}

When you want to convert \TEX\ documents into \XML\ (or
\HTML) counterparts, it makes sense to use \TEX\ itself,
because that way we have the full engine at our fingertips.

In \CONTEXT\ we use a strange but charming way to handle
this conversion: we simply typeset \XML. The typeset
document (with lots of angle brackets) is then converted
into raw text. A complication is math, but given that you
can embed \MATHML\ in your document, this can be passed on
directly.

\SomeLogo{102} \paragraaf[examples]{examples of markup}

\SomeLogo{4} \paragraaf[reference]{command reference}

In this chapter we summarize the mapping commands. You
should be aware of the fact that some of them are
experimental and so their names and even their meanings may
change.

For the moment we left out the commands that are used for
two||step remapping \MATHML\ onto \TEX. Once these have
proven to be useful, they will be discussed here too.
These two||step remapping commands convert an \XML\ stream
into a \TEX\ one that then can be processed in a for \TEX\
more comfortable way.

\subonderwerp {mapping elements}

\starttypen
\defineXMLsingular    [name] {command}
\defineXMLcommand     [name] {command}
\defineXMLgrouped     [name] {command}
\defineXMLargument    [name] {command}
\defineXMLargument    [name] {command}
\defineXMLpickup      [name] {before} {after}
\defineXMLenvironment [name] {before} {after}
\defineXMLpush        [name]
\defineXMLprocess     [name]
\defineXMLnested      [name] {before} {after}
\stoptypen

\subonderwerp{instruction processors}

\starttypen
\defineXMLprocessor [name] \command
\stoptypen

The following processing instructions are predefined:

\starttypen
<?context-command   \somecommand{to be executed} ?>
<?context-directive class key value              ?>
<?context-message   some text to show in the log ?>
\stoptypen

\subonderwerp{mapping entities}

\starttypen
\defineXMLentity [name] {command}
\stoptypen

\subonderwerp{processing files}

\starttypen
\XMLinput               filename
\processXMLfile        {filename}
\processXMLfilegrouped {filename}
\showXMLfile           {filename}
\stoptypen

\subonderwerp{embedding code}

\starttypen
\startXMLdata xml commands \stopXMLdata

\processXMLbuffer [name]

\startXMLcode [name] xml commands \stopXMLcode
\getXMLcode   [name]
\stoptypen

\subonderwerp{loading filters}

\starttypen
\useXMLfilter [list]
\stoptypen

\subonderwerp{handling attributes}

\starttypen
\mapXMLvalue {name} {key} {old} {new}

\XMLvar {name} {key} {default}
\XMLval {name} {key} {default}
\XMLpar {name} {key} {default}

\doifXMLvarelse {name} {key} {then branch} {else branch}
\doifXMLvalelse {name} {key} {then branch} {else branch}
\doifXMLparelse {name} {key} {then branch} {else branch}

\getXMLarguments {class} {key="value" pairs}

\XMLifequalelse {name} {key} {value} {then branch} {else branch}
\stoptypen

\subonderwerp{manipulating pushed elements}

\starttypen
\XMLpop     {name}
\XMLappend  {name} {text}
\XMLprepend {name} {text}
\XMLerase   {name}
\XMLassign  {name}
\XMLshow    {name}

\defXMLstring \command {name}

\doifXMLdataelse {name} {then branch} {else branch}
\stoptypen

\subonderwerp{auxiliary commands}

\starttypen
\defXMLlowerclean \command

\bXMLs
\eXMLs
\stoptypen

% \SomeChart {conversion 1} {TEX     1}
% \SomeChart {conversion 1} {CONTEXT 1}
% \SomeChart {conversion 1} {PDF     1}
%
% \SomeChart {conversion 2} {XML     2}
% \SomeChart {conversion 2} {CONTEXT 2}
% \SomeChart {conversion 2} {PDF     2}
%
% \SomeChart {conversion 3} {XML     3}
% \SomeChart {conversion 3} {TEX     3}
% \SomeChart {conversion 3} {CONTEXT 3}
% \SomeChart {conversion 3} {PDF     3}
%
% \SomeChart {conversion 4} {XML     4}
% \SomeChart {conversion 4} {XSLT    4}
% \SomeChart {conversion 4} {TEX     4}
% \SomeChart {conversion 4} {CONTEXT 4}
% \SomeChart {conversion 4} {PDF     4}
%
% \SomeChart {conversion 6} {XML     6}
% \SomeChart {conversion 6} {XSLT    6}
% \SomeChart {conversion 6} {XML     6}
% \SomeChart {conversion 6} {CONTEXT 6}
% \SomeChart {conversion 6} {PDF     6}
%
% \SomeChart {conversion 5} {TEX     5}
% \SomeChart {conversion 5} {CONTEXT 5}
% \SomeChart {conversion 5} {XML     5}

%\SomeLogo{101} \paragraaf[educational]{educational markup}
%
%Although over time, \CONTEXT\ has been used for all kind of
%applications, the primary reason for its existence is that
%we needed a tool to typeset complex educational documents.
%
%
%An important characteristic of such documents is that they
%can have many different ways of presenting information, as
%well as that they are quite demanding in terms of reusing
%components. This means that one has to provide a lot of
%structure, and since \XML\ is about structure, educational
%documents make a nice example.
%
%Here we will not so much present complete styles and \XML\
%definitions, but introduce some general purpose techniques
%for handling this kind of data.

%\SomeLogo{102} \paragraaf[examination]{examination markup}
%
%Every Dutch school kid at some moment has to pass an
%official examn. When in the mid nineties we were in search
%of a good example for interactive documents, we decided to
%take the examns as example source code. That way we could
%make useful documents, free of copyright, and demonstrating
%a broad range of typographical challenges.
%
%The examns were either rekeyed from print, or edited from
%rough scanned and converted documents. \unknown\

\pagina \stelkopnummerin[paragraaf][+1]

\definieeroverlay[vooruit][\overlaybutton{eerstepagina}]

\startstandaardopmaak

\stelinteractiein[letter=]

\steluitlijnenin[midden]

\startkleur[backgroundcolor]

\vskip24pt

\leavevmode\externalfigure
  [\bufferprefix mprun.304][hoogte=.75\teksthoogte]

\vfill

\AuthorFont \stelinterliniein

PRAGMA Advanced Document Engineering

\vskip6pt

\naar{www.pragma-ade.com}[url(http://www.pragma-ade.com)]

\vskip24pt

\stopkleur

\stopstandaardopmaak

\stoptekst
